\section{Four missing witnesses and a smaller list of cases}\label{sec:four}
This section proves the full lower bound when four missing witnesses increase the cover number by four. The argument begins with a product-count obstruction, then treats the possible ways the higher-degree contributions enter the quadratic circuit. Let $H$ be the quadratic support graph, and let $C_2$ be the circuit obtained by minimum-degree extraction followed by pruning. Define
\[
 d=\tau(G)-\tau(H),\qquad D=|C|-|C_2|,\qquad
 e=|C_2|-(|E(H)|+\tau(H)-1).
\]
Thus $d$ is the increase in cover number, $D$ counts deleted gates, and $e$ is the excess size of the quadratic circuit. In the four-witness case, the existing reduction leaves precisely the rows of \citet[Corollary~5.34]{Companion}. We remove the pair $(D,e)=(6,0)$ from its $d=4$ row.

Here is the extraction used throughout the argument. The minimum degree
$\delta$ is one at a variable, is the sum of the input degrees at a
product, and is their minimum at an addition. Starting from the output,
select both inputs at a product or an equal-degree addition, and only the
lower-degree input at an unequal-degree addition. Bypass each selected
unequal-degree addition and discard unused gates. The resulting circuit
$C_2$, called the \emph{quadratic frame}, computes the degree-two part of
$P$; positivity prevents cancellation. A \emph{private} gate is absent
from this retained circuit. An unequal-degree addition is a
\emph{cut}; its lower- and higher-degree operands are its low and high
inputs. A cut is \emph{visible} if its own node belongs to the selected
output subgraph before bypassing. Reuse of its low input elsewhere does
not make the cut visible.

For the Boolean arguments, a function is \emph{safe} for $G$ if every
assignment on which it is true contains an edge of $G$. It
\emph{witnesses} an edge if it is true on the assignment consisting of
that edge's two endpoints. A positive clause is a disjunction of input
variables; the empty clause is zero. An affine positive clause may also
have a constant term, in which case a positive constant makes it one.
A conjunctive normal form (CNF) is a conjunction of such clauses.
These constants are permitted in the relaxed Boolean templates and in
restricted circuits, even though the original circuit has only variable
inputs. Positive arithmetic linear forms give these clauses on Boolean
supports.


The following finite lemma is proved by an explicit propositional certificate. It is used only as an obstruction on a larger class: the template deliberately omits the exact-frame, optimality and pruning restrictions that a circuit would satisfy. Thus it is not a counterexample search, and its unsatisfiability cannot discard a legitimate circuit.

\begin{lemma}[A three-product Boolean template]\label{lem:threeprivatefinite}
For $i=0,1,2$, let
\[
 g_i=A_i(x,g_0,\ldots,g_{i-1})\wedge B_i(x,g_0,\ldots,g_{i-1}),
\]
where $x=(x_0,\ldots,x_7)$ and $A_i,B_i$ are arbitrary disjunctions of subsets of the displayed arguments; the empty disjunction is allowed.
Let $L_i$ be arbitrary affine positive clauses in $x$, including zero and one, and let $c_{ij}\in\{0,1\}$.
Then
\[
 \bigvee_i(g_i\wedge L_i)\ \vee\!
 \bigvee_{i<j:c_{ij}=1}(g_i\wedge g_j)
 \ne\bigvee_{a=0}^3(x_{2a}\wedge x_{2a+1}).
\]
\end{lemma}
\begin{proof}[Proof by finite certificate]
Use one selector for each possible argument in each input clause, each possible variable or constant in each $L_i$, and each cross term. There are $54+27+3=84$ selectors. For fixed selectors the function is monotone, so equality with the matching is equivalent to truth on its four edge assignments and falsity on its sixteen maximal independent endpoint transversals.
The program \emph{Private-product obstruction} expresses evaluations on these twenty assignments by exact Tseitin equivalences: a selected argument is the conjunction of its selector and its value, each input is their disjunction, and each product and output gate is encoded in both directions.
The resulting CNF has $2{,}085$ variables and $7{,}421$ clauses.

Its DRUP proof contains $43{,}308$ additions and $41{,}308$ deletions. An independent checker using only the C++ standard library verifies every addition by reverse unit propagation and derives the empty clause. Deletions are ignored, soundly retaining only original or already implied clauses. The check took $15.18$ seconds. The encoder, proof and standard-library checker are available through the workbench. Reverse unit propagation means that adjoining the negation of the proposed clause and propagating forced literals yields a contradiction. Therefore each accepted clause follows from the earlier clauses; induction ends in a proof of inconsistency. The complete check is a propositional certificate.
\end{proof}

\begin{proposition}[Four missing induced edges require four deleted products]\label{prop:fourprivateproducts}
Let a positive variable-input binary circuit compute $P\in\Sup(G)$ with nonempty quadratic support $H$. Suppose eight vertices induce precisely four disjoint edges of $G$ and no edge of $H$. Minimum-degree extraction and pruning delete at least four product gates. There is no bound on deleted additions and no quadratic-optimality or cover-number hypothesis.
\end{proposition}
\begin{proof}
Suppose at most three products are deleted. Apply the full-carrier erasure from \citet[proof of Lemma~5.26, first three paragraphs]{Companion}: protect the original minimum-degree-two nodes with addition-only output paths, and erase their off-carrier uses while their \emph{whole current values} still reach the output additively. The erased multiples are Boolean-absorbed. This preserves the exact extracted quadratic circuit, the Boolean output and the bound of three deleted products. It does not replace a carrier by its quadratic part.

No carrier now feeds its complement. Every product in that complement is deleted, because a quadratic product that survives extraction has an addition-only output path and is itself protected. Order those private products as $g_0,g_1,g_2$, padding with unused zero expressions if necessary. Induction through the private addition gates shows that each input to product $i$ is a nonnegative linear combination of original variables and earlier $g_j$. Booleanly this is exactly the arbitrary input clause permitted by Lemma~\ref{lem:threeprivatefinite}; arbitrary sharing is included.

The operands of a quadratic-frame product have the form
\[
 A+\sum_i a_i g_i,\qquad B+\sum_i b_i g_i,
\]
with $A,B$ linear forms and all coefficients nonnegative. Expand these quadratic-frame products, leaving the private products unexpanded. The protected output network uses additions alone, so summing gives
\[
 P_2+\sum_i g_iL_i+\sum_{i\leq j}c_{ij}g_ig_j,
\]
where $L_i$ are affine nonnegative linear forms in the original variables. Direct private-product contributions give their constant terms. No input variable can reach the output additively because the output has minimum degree two. The input-only contribution is precisely $P_2$: the quadratic-frame products multiply their original degree-one parts, and each private product has minimum degree at least two.

Set every variable outside the specified eight vertices to zero. Then $P_2$ vanishes, and the Boolean output is exactly their four-edge matching. Boolean idempotence absorbs $g_i^2$ into the constant part of $L_i$, and positive coefficients become presence bits. The result is the forbidden template of Lemma~\ref{lem:threeprivatefinite}, a contradiction.
\end{proof}

\begin{corollary}[The optimal six-deletion case]
Among the four-missing-witness alternatives, $(d,D,e)=(4,6,0)$ is impossible. The remaining necessary pairs are
\[
\begin{array}{c|l}
d&(D,e)\\\hline
2&(5,0)\\
3&(5,0),(5,1),(6,0)\\
4&(6,1),(7,0).
\end{array}
\]
\end{corollary}
\begin{proof}
For $d=4$, the missing edges form an induced matching with no old edge on their endpoints. Full-carrier erasure preserves the optimal quadratic frame, and \citet[Lemma~5.33]{Companion} forces all six deleted gates to remain. Proposition~\ref{prop:fourprivateproducts} gives at least four private products. At least one deleted addition is an unequal-degree cut, since otherwise the useful circuit is homogeneous and its degree-two output has no missing witness. Thus the product/addition profile is $(4,2)$ or $(5,1)$.

The four-product case is excluded by \citet[Lemma~5.35]{Companion}. In the five-product case there is a unique cut. Its high input is a connected multiplication DAG with at most five products and hence at most six distinct old linear boundary inputs: ten operand occurrences include at least four internal uses. Its scalar or linear direct quadratic-frame context adds at most one factor. After restriction and Boolean idempotence the matching would therefore be a conjunction of at most seven affine clauses.
Such a representation requires at least sixteen clauses: every nonconstant clause meets each of the four edge assignments and rejects at most one independent endpoint transversal, whereas all sixteen transversals must be rejected. This contradiction completes the proof.
\end{proof}

The same argument leaves only four products and two additions when $(d,D,e)=(4,6,1)$, and only the product/addition profiles $(4,3)$ and $(5,2)$ when $(d,D,e)=(4,7,0)$. These restrictions do not assert the existence of any remaining circuit. The full four-missing-witness theorem remains open here.
